\section{Application to persistence modules and parameter extension} \label{Application to persistence modules and parameter extension}
Persistence modules provide a natural connection between the diagram categories considered above and topological data analysis. Let $k$ be a field. We denote by $Vect_k$ the category of $k$-vector spaces, and by $vec_k$ its full subcategory of finite-dimensional $k$-vector spaces. A persistence module indexed by a poset $P$ is a functor $P \longrightarrow Vect_k$, and it is called pointwise finite-dimensional when all its values are finite-dimensional \cite[Section 2.3]{Botnan}. Thus, for a finite poset $P$, the diagram category $vec_k^P$ is the category of pointwise finite-dimensional persistence modules indexed by $P$.

Since $vec_k$ is an abelian category, if two finite posets $P$  and $Q$ are universally derived equivalent, then, for every field $k$, \[D(vec_k^P) \simeq D(vec_k^Q). \] Hence every universal derived equivalence between finite posets obtained above yields, in particular, a derived equivalence between the corresponding categories of pointwise finite-dimensional persistence modules.

For $n \geq 1$, let $I_n$ denote the $n$-element chain. 
When $P$ is a product of finite chains, this recovers the usual discrete multiparameter setting: a product $I_{n_1} \times \dots \times I_{n_d}$ is a finite $d$-dimensional parameter grid \cite[Sections 2.2-2.3]{Botnan}. More generally, arbitrary finite posets allow parameter spaces with branching or nonstandard orientations.

%\subsection{Parameter extension}

The following result is a special case of Ladkani \cite[Lemma 2.9]{ladkani}. We include a short argument for completeness.

\begin{proposition}%[Stability under parameter extension] 
\label{Stability under parameter extension}
    Let $X$ and $Y$ be finite universally derived equivalent posets, and let $R$ be any finite poset. Then
    \[X \times R \,\,\, \text{ and } \,\,\, Y \times R\] are universally derived equivalent.
\end{proposition}

\begin{proof}
    Let $\mathcal{A}$ be an arbitrary abelian category. Since $R$ is finite, the diagram category $\mathcal{A}^R$ is again abelian. Moreover, there are natural isomorphisms of categories 
    \[\mathcal{A}^{X \times R} \cong (\mathcal{A}^R)^X \,\,\, \text{ and }\,\,\, \mathcal{A}^{Y \times R} \cong (\mathcal{A}^R)^Y.\] Since $X$ and $Y$ are universally derived equivalent, we may apply the defining equivalence to the abelian category $\mathcal{A}^R$. Hence \[D((\mathcal{A}^R)^X) \simeq D((\mathcal{A}^R)^Y).\] Using the natural identification above, we obtain \[D(\mathcal{A}^{X \times R}) \simeq D(\mathcal{A}^{Y \times R}).\] Since $\mathcal{A}$ was arbitrary, $X \times R$ and $Y \times R$ are universally derived equivalent. 
\end{proof}
\begin{corollary} %[Parameter extension for admissible cuts] 
\label{Parameter extension for admissible cuts}
    Let $P=A \sqcup B$ be an admissible cut and let $P^-$ be its associated negative poset. Then, for every finite poset $R$, 
    \[P \times R \,\,\, \text{ and } \,\,\, P^- \times R\] are universally derived equivalent.
\end{corollary}

\begin{proof}
    By Corollary \ref{The associated universally derived equivalent poset}, $P$ and $P^-$ are universally derived equivalent. The conclusion then follows directly from Proposition \ref{Stability under parameter extension}.
\end{proof}

In the persistence setting, Corollary \ref{Parameter extension for admissible cuts} allows one to adjoin further discrete parameters to any universally derived equivalent pair produced by an admissible cut. In particular, taking
\[R=I_{n_1} \times \dots \times I_{n_d}\] yields derived equivalences between categories of persistence modules obtained by adjoining $d$ ordinary discrete persistence parameters.

\begin{example} \label{example Tamari TDA}
    Consider the admissible cut of the Tamari lattice $T_3$ from Example \ref{T3/D5}, and let $T_3^-$ be its associated negative poset. By Corollary \ref{Parameter extension for admissible cuts}, for every finite poset $R$,
    \[T_3 \times R \,\,\, \text{ and } \,\,\, T_3^- \times R\] are universally derived equivalent.

    In particular, taking $R=I_n$ yields, for every $n \geq 1$, a pair of universally derived equivalent posets with $5n$ elements. Thus the equivalence of Example \ref{T3/D5} gives rise to an infinite family of universally derived equivalent pairs. 

    More generally, if
    \[R=I_{n_1} \times \dots \times I_{n_d},\] then, for every field $k$,
    \[D(vec_k^{T_3 \times R}) \simeq D(vec_k^{T_3^- \times R}).\] Thus, the equivalence of Example \ref{T3/D5} extends to parameter spaces obtained by adjoining any finite number of ordinary discrete parameters.

    When $R$ is a product of nonempty finite chains, these two parameter posets are still non-isomorphic. Indeed, $T_3 \times R$ has a unique maximal element, whereas $T_3^- \times R$ has two maximal elements. Thus the construction yields infinite families of non-isomorphic parameter posets with equivalent derived categories of pointwise finite-dimensional persistence modules.
\end{example}

\begin{corollary}%[Parameter extension for tree orientations] 
\label{Parameter extension for tree orientation}
    Let $T$ be a finite tree, let $P,Q \in \operatorname{Or}(T)$, and let $R$ be any finite poset. Then
    \[P \times R \,\,\, \text{ and } Q \times R\] are universally derived equivalent.
\end{corollary}

\begin{proof}
    By Theorem \ref{tree-exhaustiveness}, $P$ and $Q$ are universally derived equivalent. The conclusion follows from Proposition \ref{Stability under parameter extension}.
\end{proof}


In particular, if
\[R=I_{n_1} \times \dots \times I_{n_d},\] then the posets $P \times R$, as $P$ ranges over $\operatorname{Or(T)}$, yield derived equivalent categories of pointwise finite-dimensional persistence modules. Thus, after adjoining any finite number of ordinary discrete parameters, the orientation of the tree coordinate is not detected by the derived category.

The following example illustrates Corollary \ref{Parameter extension for tree orientation} in the simplest nontrivial two-parameter case.
\begin{example}
    Let $I_3$ and $Z_3$ be the following two orientations of the three-vertex tree, where $I_3$ is the three-element chain, and let $I_2$ denote the two-element chain. Set
    \[G \coloneq I_3 \times I_2, \,\,\,\,\, H \coloneq Z_3 \times I_2.\]
   
% https://q.uiver.app/#q=WzAsOCxbMiwxLCJ4XzEiXSxbMCwyLCJ4XzAiXSxbNCwyLCJ4XzIiXSxbOCwyLCJ6XzAiXSxbMTAsMSwiel8xIl0sWzEyLDIsInpfMiJdLFsyLDAsIkNfMyJdLFsxMCwwLCJaXzMiXSxbMSwwXSxbMCwyXSxbMyw0XSxbNSw0XV0=
\[\begin{tikzcd}[sep=small]
	&& {I_3} &&&&&&&& {Z_3} && \\
	&& {x_1} &&&&&&&& {z_1} \\
	{x_0} &&&& {x_2} &&&& {z_0} &&&& {z_2}
	\arrow[from=2-3, to=3-5]
	\arrow[from=3-1, to=2-3]
	\arrow[from=3-9, to=2-11]
	\arrow[from=3-13, to=2-11]
\end{tikzcd}\]
\\

% https://q.uiver.app/#q=WzAsMTQsWzIsMSwieF8xIl0sWzAsMSwieF8wIl0sWzQsMSwieF8yIl0sWzgsMSwiel8wIl0sWzEwLDEsInpfMSJdLFsxMiwxLCJ6XzIiXSxbMiwwLCJHPUNfMyBcXHRpbWVzIElfMiJdLFsxMCwwLCJIPVpfMyBcXHRpbWVzIElfMiJdLFswLDMsInhfMCciXSxbMiwzLCJ4XzEnIl0sWzQsMywieF8yJyJdLFs4LDMsInpfMCciXSxbMTAsMywiel8xJyJdLFsxMiwzLCJ6XzInIl0sWzEsMF0sWzAsMl0sWzMsNF0sWzUsNF0sWzEsOF0sWzgsOV0sWzAsOV0sWzksMTBdLFsyLDEwXSxbMywxMV0sWzQsMTJdLFsxMSwxMl0sWzUsMTNdLFsxMywxMl1d
\[\begin{tikzcd}[column sep=small,row sep=scriptsize]
	&& {G=I_3 \times I_2} &&&&&&&& {H=Z_3 \times I_2} && \\
	{x_0} && {x_1} && {x_2} &&&& {z_0} && {z_1} && {z_2} \\
	\\
	{x_0'} && {x_1'} && {x_2'} &&&& {z_0'} && {z_1'} && {z_2'}
	\arrow[from=2-1, to=2-3]
	\arrow[from=2-1, to=4-1]
	\arrow[from=2-3, to=2-5]
	\arrow[from=2-3, to=4-3]
	\arrow[from=2-5, to=4-5]
	\arrow[from=2-9, to=2-11]
	\arrow[from=2-9, to=4-9]
	\arrow[from=2-11, to=4-11]
	\arrow[from=2-13, to=2-11]
	\arrow[from=2-13, to=4-13]
	\arrow[from=4-1, to=4-3]
	\arrow[from=4-3, to=4-5]
	\arrow[from=4-9, to=4-11]
	\arrow[from=4-13, to=4-11]
\end{tikzcd}\]
 By Theorem \ref{tree-exhaustiveness}, the posets $I_3$ and $Z_3$ are universally derived equivalent, since they are two orientations of the same finite tree. Therefore, by Proposition \ref{Stability under parameter extension}, the product posets $G$ and $H$ are universally derived equivalent.

    The poset $G$ is the usual $3 \times 2$ rectangular grid, which is a standard finite parameter space for discrete two-parameter persistence \cite[Sections 2.2-2.3]{Botnan}. By contrast, $H$ combines one ordinary discrete parameter with one zigzag-shaped parameter. Thus $G$ and $H$ provide an example of two non-isomorphic parameter posets whose categories of pointwise finite-dimensional persistence modules have equivalent derived categories. 

    To see that $G$ and $H$ are not isomorphic, it is enough to compare their minimal elements: $G$ has a unique minimal element, whereas $H$ has two minimal elements. This shows that universal derived equivalence may relate a standard multiparameter grid to a parameter poset with a nonstandard orientation.

\end{example}

\section*{Acknowledgments}
I would like to thank Jorge Vitória for his invaluable guidance, many helpful discussions, and numerous suggestions since the early stages of this work. I am also grateful to Eero Hyry for reading an earlier version of the manuscript and for useful comments.